\section{证据约束的抽取与真值表}
\label{sec:extraction}

\begin{figure}[t]
  \centering
  \figslot{fig2_extraction}{5.6cm}
  \caption{证据约束的双通道抽取。规则通道从白名单表格中按科目匹配取数；模型通道（分析师会话按页读 PDF）给出原文写法与页码。程序逐字回查原文、合并两通道、做勾稽，任何拒收、冲突、勾稽失败或核心缺失都会让流程回到抽取（最多三轮），通过后才生成真值表。}
  \label{fig:extraction}
\end{figure}

抽取是整个系统的地基：真值表中的一个错数会沿着分析、估值、文稿一路传播。图 \ref{fig:extraction} 给出了抽取流程。本节按顺序介绍页定位、面向真实年报版式的解析、两条抽取通道、合并规则与校验闸门。

\subsection{页定位}
一份 A 股年报通常有 200 页以上，三大报表和主要会计数据只占其中十几页。\code{locate} 扫描全文，给出建议阅读页、各页的单位声明、母公司报表所在页（需要回避），并标出没有文字层的扫描页。扫描页目前直接拒收，不做光学字符识别：与其引入一个错误率未知的识别环节，不如明确告知输入不合格。

\subsection{真实年报版式的解析}
\label{sec:parsing}
在自造样例上，表格解析几乎没有难度；换到真实年报，规则通道起初只能取到 314 个核验值中的 128 个。逐一排查后，我们归纳出四类版式问题并分别处理：

\begin{itemize}
  \item \textbf{稀疏网格。}部分年报（如宁德时代）的表格在 PDF 中被切成远多于实际列数的网格，科目与数字落在不同的空列里。\code{compact\_columns} 合并稀疏列：从左到右把相邻列归为一组，只要组内任一行两列都有值就另起一组；规整的表格没有全空列，原样返回。
  \item \textbf{跨页续表。}资产负债表等长表跨页时，续页没有表头，规则通道无法判断各列对应的期间。\code{link\_continuation} 判断页首无标题的表格是否为上一页表格的续表：若是，则沿用首页表头的期间列，并记录 \code{continued\_from}。我们只记录续表关系、不合并两页的行，因此每一行的页码仍是它实际所在的页，来源定位不失真。
  \item \textbf{科目折行。}“归属于上市公司股东的净利润”这类长科目会折成两行，数字只在第一行。\code{merge\_wrapped\_labels} 把只有标签、没有数值、且满足折行特征（上一行括号未闭合、上一行写满列宽、或本行以不会作科目开头的字开头）的行接回上一行；以冒号结尾的小节标题不算折行。
  \item \textbf{同页多单位。}同一页可能先后出现“单位：元”与“单位：万元”。解析器记录各“单位：”声明的纵向位置，每张表取其上方最近的声明；表格标题只在上一张表格底部以下、最多向上六行范围内查找。
\end{itemize}

\subsection{规则通道}
规则通道只从白名单表格取数：主要会计数据、分季度主要财务指标、非经常性损益、三大报表及其补充资料。真实年报附注里有大量与报表科目同名但口径不同的明细表（受限资产、费用明细、关联方往来），混入会造成冲突，因此被排除。具体规则包括：标题含“母公司”的表格不取（只取合并报表）；首行含“关联方”的表格不取；表头含“比例”“比重”“占”的结构比列不当作期间列；同表既有“营业总收入”又有“营业收入”时（含金融子公司利息收入等情形），营业收入取后者；现金流量表补充资料中的调节项（如“财务费用（收益以‘－’号填列）”）不当作报表科目本身。财报自带的“本年比上年增减”列也被保留，用于与程序计算的同比交叉核验；比率类指标（如净资产收益率）的增减按百分点核对，而不是相对增速。

\subsection{模型通道与证据校验}
分析师会话按 \code{locate} 的建议逐页阅读，写出 \file{facts.json}：每条事实给出指标、期间、原文写法与页码。

模型给出的每个数字都要经过证据校验（\code{find\_evidence}）：设 $P$ 为该条所写的页码，$\nu(r)$ 为原文写法 $r$ 去掉千分位、括号、百分号与正负号后的数字串，则要求该页至少有一个片段（表格优先）的文本按\textbf{数字边界}包含 $\nu(r)$，
\begin{equation}
  \exists\, c_k:\ \mathrm{page}(c_k)=P\ \wedge\ \nu(r)\ \text{以完整数字的形式出现在}\ \mathrm{text}(c_k)\ \text{中},
\end{equation}
即前后不能紧邻其他数字或小数点，避免“1.23”命中“11.234”这类部分匹配。此外，期间格式不合法、口径不是合并报表、单位不合法、缺页码的条目同样拒收。否则该条被拒收并记录原因。这条规则简单到近乎粗暴，却有效地排除了模型“算出来的”或“记忆中的”数字：模型可以抄错位置，但很难凭空写出一个恰好出现在所引页上的数字。

\subsection{合并：不静默取舍}
算法 \ref{alg:merge} 给出合并规则。同一通道内多处取到同一 $(m,p)$ 时，取值相同则互为佐证，不同则生成高严重度问题；两通道之间取值相同标记为 \code{rule+llm}，不同则生成中严重度问题，并注明暂时保留的是模型读原文的结果。无论保留哪一方，问题记录都会进入抽取闸门。

\begin{algorithm}[t]
\caption{双通道合并（\code{merge}）}
\label{alg:merge}
\small
\begin{algorithmic}[1]
\Require 规则通道事实 $F_R$，模型通道事实 $F_L$，冲突时的保留方 $\pi\in\{\text{rule},\text{llm}\}$
\State $M\gets\varnothing$，$I\gets\varnothing$
\For{$f\in F_R$}
  \If{$(m,p)\notin M$} $M[(m,p)]\gets f$
  \ElsIf{$\mathrm{close}(M[(m,p)].v, f.v)$} 记 $f.\sigma$ 为佐证
  \Else{} $I\gets I\cup\{\text{高：同一指标多处取值不一致}\}$
  \EndIf
\EndFor
\For{$f\in F_L$}
  \If{$(m,p)\notin M$} $M[(m,p)]\gets f$
  \ElsIf{$\mathrm{close}(M[(m,p)].v, f.v)$} $M[(m,p)].\text{method}\gets\text{rule+llm}$
  \Else{} $I\gets I\cup\{\text{中：规则与模型不一致（注明保留方）}\}$；\textbf{if} $\pi=\text{llm}$ \textbf{then} $M[(m,p)]\gets f$
  \EndIf
\EndFor
\State \Return $M, I$ \Comment{$I$ 非空即阻断估值}
\end{algorithmic}
\end{algorithm}

\subsection{校验闸门与勾稽}
\code{verify-facts} 把合并后的事实送入分析模块，再汇总四类阻断问题：拒收（证据校验失败）、冲突（算法 \ref{alg:merge} 的问题记录）、勾稽失败与核心指标缺失。勾稽包括：
\begin{itemize}
  \item 资产总计 $=$ 负债和所有者权益总计；负债和所有者权益总计 $=$ 负债合计 $+$ 所有者权益合计；所有者权益合计 $=$ 归母权益 $+$ 少数股东权益（相对容差 $10^{-6}$，且不小于 1 元）；
  \item 非经常性损益合计 $=$ 归母净利润 $-$ 扣非归母净利润（仅当非经常性损益为财报列示值时核对）；
  \item 财报列示的同比增减与程序计算的同比在展示精度内一致；
  \item 单季拆分：$Q_2=H_1-Q_1$，$Q_3=Q_{1\text{-}3}-H_1$，$Q_4=\text{全年}-Q_{1\text{-}3}$；当某个单季数既有财报列示值、又能由累计数相减得出时，二者必须一致。程序记录每个单元格依赖的列示单元格集合，避免“自己和自己比”。
\end{itemize}
有任何阻断问题时命令返回非零，模型须回到原文核对修正，最多三轮。只有全部通过，才会生成 \file{facts.verified.json} 并进入估值。

\subsection{经营数据}
估值需要的分产品量价、区域收入、产能、分红等数据不在三大报表里。模型按分析框架中的驱动清单把相关表格逐格抄进 \file{extras.json}，\code{verify-extra} 逐格回查原文；任何派生比率（吨价、增速、差值）不允许模型自己写，只能声明为 \code{ratios}，由程序计算，并支持别名以便正文引用。
